Os avanços na interação humano-computador (HCI) dependem de sensores acessíveis e precisos para rastreamento de movimentos. \citep{guzsvinecz2019suitability} revisaram a adequação dos sensores Kinect e Leap Motion Controller (LMC) em aplicações multidisciplinares, como reabilitação médica e educação.

Os dispositivos foram comparados quanto à precisão, custo e algoritmos de reconhecimento de gestos. Estudos anteriores demonstraram a utilidade do Kinect para rastreamento corporal completo \citep{redaelli2021comparison}, enquanto o LMC destacou-se no rastreamento de movimentos das mãos em tarefas que exigem alta sensibilidade \citep{weichert2013analysis}. Embora menos precisos que sensores de alta gama, esses dispositivos oferecem uma solução acessível para tarefas específicas \citep{guzsvinecz2019suitability}, sendo especialmente adequados para aplicações educacionais e médicas que requerem interatividade com custos reduzidos.

O uso de interfaces naturais para controle remoto de dispositivos robóticos tem se tornado um tema relevante na pesquisa em interação humano-máquina. \citep{moura2016estudos} exploraram o potencial do sensor Kinect como alternativa para dispositivos tradicionais, destacando sua capacidade de rastrear movimentos humanos com precisão e replicá-los em um braço robótico construído com peças LEGO.
O estudo demonstrou que a interface proposta facilita a interação, sendo considerada de fácil utilização por voluntários, apesar de limitações em termos de velocidade e flexibilidade de movimentos. A pesquisa reafirma a importância do desenvolvimento de interfaces acessíveis e intuitivas, especialmente para aplicações educacionais e de inclusão social.


